Borcea--Br\"and\'en posed the relevant representation and endpoint problems in the AIM problem list \cite{AIMStableProblems}.  Let $S=x_1+\cdots+x_5$ and
\begin{equation}\label{eq:stable-p}
p(x_1,\ldots,x_5)=\prod_{i=1}^5\left(x_i+5\sum_{j\ne i}x_j\right)=\prod_{i=1}^5(5S-4x_i).
\end{equation}
Every factor has strictly positive coefficients, hence has positive imaginary part whenever all variables do; therefore $p$ is real stable.  It is symmetric, homogeneous of degree five in five variables, and has strictly positive monomial coefficients.

\begin{theorem}[\statusfalse: Stable Permanent-On-Top, AIM Problem 38]\label{thm:pot}
The polynomial \eqref{eq:stable-p} satisfies all hypotheses of Borcea--Br\"and\'en Problem 36 but violates
\[
a_\lambda\le f^\lambda a_{(1^5)}.
\]
\end{theorem}
\begin{proof}
Its exact Schur expansion is
\begin{align*}
p={}&625s_{(5)}+4500s_{(4,1)}+7125s_{(3,2)}+8150s_{(3,1,1)}\\
&+7525s_{(2,2,1)}+6580s_{(2,1,1,1)}+1281s_{(1^5)}.
\end{align*}
For $\lambda=(3,2)$, the hook-length formula gives $f^{(3,2)}=5$, whereas
\[
7125>5\cdot1281=6405.
\]
The same upper endpoint fails for $(3,1,1)$, $(2,2,1)$, and $(2,1,1,1)$.  The source package independently reconstructs the monomial expansion and inverts the Kostka matrix over $\mathbb Z$.
\end{proof}

The polynomial also has a diagonal positive-definite mixed-determinant representation.  For $j=1,\ldots,5$, let
\[
D_j=\diag(5,\ldots,5,\underset{j}{1},5,\ldots,5).
\]
Then the multivariate Johnson polynomial satisfies
\[
\eta(x_1D_1,\ldots,x_5D_5)=p(x_1,\ldots,x_5).
\]

\begin{theorem}[\statusfalse: common-matrix representation, AIM Problem 35]\label{thm:common-A}
There is no $5\times5$ Hermitian positive semidefinite matrix $A$ such that
\[
p(x_1,\ldots,x_5)=\eta(x_1A,\ldots,x_5A).
\]
\end{theorem}
\begin{proof}
Such a representation would identify the Schur coefficients, up to one common positive normalization, with irreducible immanants of $A$ indexed by conjugate partitions.  Immanant Permanent-On-Top is known in size five, indeed in every size at most thirteen \cite{Wanless2022}.  It would therefore force
$a_\lambda\le f^\lambda a_{(1^5)}$ for every $\lambda\vdash5$, contradicting Theorem~\ref{thm:pot}.
\end{proof}

\begin{remark}[Scope]
The example is Schur-positive, and its lower endpoint inequalities hold, so it refutes neither Problem 36 nor Problem 37; those require the positive-definite pencil of Theorem~\ref{thm:aim36} below.  The full independent audit is included as \path{counterexamples/stable-schur/artifacts/}.
\end{remark}
